\section{Tuần 2: Trí tuệ Nhân tạo, Học máy \& Học sâu trong Khoa học Dữ liệu Thị giác}
\label{sec:week2}

\subsection{Bản chất của Trí tuệ Tự nhiên, Trí tuệ Nhân tạo và Khái niệm Thông minh}

Trong khoa học nhận thức và khoa học máy tính, việc phân định rõ các dạng thức trí tuệ là tiền đề quan trọng để xây dựng và phát triển các hệ thống tính toán:

\textbf{Trí tuệ Tự nhiên (Natural Intelligence - NI):} Là năng lực sinh học bẩm sinh của con người và các loài sinh vật bậc cao, được tôi luyện qua hàng triệu năm tiến hóa. Trí tuệ tự nhiên tích hợp toàn diện các phương thức: thu nhận cảm giác đa giác quan, tự ý thức (Self-awareness), tư duy trừu tượng, sáng tạo nghệ thuật, suy luận quan hệ nhân quả và khả năng thích ứng linh hoạt trước những biến động bất định của môi trường sống, đi kèm với đời sống cảm xúc và các chuẩn mực đạo đức xã hội.

\textbf{Trí tuệ Nhân tạo (Artificial Intelligence - AI):} Là một phân ngành của khoa học máy tính nghiên cứu các nguyên lý toán học, cấu trúc dữ liệu và giải thuật nhằm kiến tạo các hệ thống tính toán (phần mềm và phần cứng) có khả năng mô phỏng các khía cạnh hành vi thông minh của con người. Các hành vi này bao gồm năng lực học hỏi từ dữ liệu (Learning), suy luận logic (Reasoning), nhận thức thị giác (Visual Perception), xử lý ngôn ngữ tự nhiên (NLP) và lập kế hoạch ra quyết định (Decision Making).

Bảng~\ref{tab:comp_ni_ai} đối chiếu các khía cạnh nền tảng giữa Trí tuệ Tự nhiên và Trí tuệ Nhân tạo.

\begin{table}[ht]
\centering
\small
\caption{Đối chiếu nền tảng giữa Trí tuệ Tự nhiên (NI) và Trí tuệ Nhân tạo (AI)}
\label{tab:comp_ni_ai}
\renewcommand{\arraystretch}{1.3}
\begin{tabularx}{\textwidth}{|p{3.0cm}|X|X|}
\hline
\textbf{Tiêu chí đối chiếu} & \textbf{Trí tuệ Tự nhiên (NI)} & \textbf{Trí tuệ Nhân tạo (AI)} \\
\hline
\textbf{1. Cơ chất vật lý} & Mạng nơ-ron sinh học hữu cơ ($\sim 86$ tỷ nơ-ron, $10^{14}$ khớp thần kinh synapse) hoạt động thông qua tín hiệu điện sinh học và hóa sinh. & Phần cứng bán dẫn vô cơ (Silicon, GPU/TPU) xử lý thông tin thông qua các xung điện nhị phân (bit $0$ và $1$). \\
\hline
\textbf{2. Nguồn gốc \& Mục đích} & Hình thành qua quá trình chọn lọc tự nhiên và tiến hóa sinh học; mục đích cốt lõi là sinh tồn, thích nghi sinh thái và duy trì nòi giống. & Do con người thiết kế và lập trình; mục đích là công cụ tự động hóa, giải quyết bài toán kỹ thuật và phục vụ lợi ích kinh tế - xã hội. \\
\hline
\textbf{3. Cơ chế học tập} & Học hiệu quả từ số ít mẫu (Few-shot), tiếp thu qua trải nghiệm tương tác với thế giới thực và duy trì quá trình học liên tục suốt đời. & Dựa trên tối ưu hóa thống kê trên các tập dữ liệu số khổng lồ; dễ gặp hiện tượng quên tri thức cũ (Catastrophic Forgetting) khi đổi miền dữ liệu. \\
\hline
\textbf{4. Bản chất tư duy} & Suy luận quan hệ nhân quả (Causal Reasoning), trực giác trừu tượng hóa và hiểu rõ bản chất nguyên nhân (\textbf{``Tại sao''}). & So khớp mẫu và tính toán tương quan xác suất ($P(Y \mid X)$), xấp xỉ hàm toán học trả lời câu hỏi dự đoán (\textbf{``Cái gì''}). \\
\hline
\textbf{5. Ý thức \& Cảm xúc} & Có ý thức tự thân, cảm xúc thực và gắn liền với sự thấu cảm cùng trách nhiệm đạo đức. & Hoạt động dựa trên xử lý tín hiệu và tính toán số học thuần túy; phản hồi cảm xúc chỉ là kết quả mô phỏng bề mặt qua văn bản hoặc hình ảnh. \\
\hline
\textbf{6. Hiệu suất năng lượng} & Tối ưu hóa cao: toàn bộ não bộ con người chỉ tiêu thụ mức công suất sinh học ước tính khoảng $20\text{ W}$. & Tiêu hao năng lượng lớn: các trung tâm dữ liệu huấn luyện mô hình nền tảng tiêu thụ từ hàng kilowatt đến hàng megawatt điện. \\
\hline
\textbf{7. Phạm vi năng lực} & Trí tuệ tổng quát (General Intelligence): con người có thể học đồng thời nhiều kỹ năng khác nhau và chuyển giao tri thức linh hoạt. & Trí tuệ chuyên biệt (Narrow AI): thực thi vượt trội trong tác vụ hẹp được huấn luyện (cờ vây, nhận diện khuôn mặt) nhưng khó tổng quát hóa sang miền mới. \\
\hline
\end{tabularx}
\end{table}

Xét về mặt vận hành, mọi hệ thống thông minh đều thực thi một vòng lặp nhận thức khép kín gồm 3 giai đoạn liên hoàn:
\begin{enumerate}
    \item \textbf{Giai đoạn 1: Nhận biết (Perception):} Tiếp nhận các kích thích và tín hiệu vật lý từ môi trường bên ngoài. Ở con người, đó là hệ thống thụ cảm đa giác quan (thị giác, thính giác, xúc giác) với cơ chế chú ý tự nhiên chọn lọc thông tin. Ở AI, đó là các cảm biến quang học, LiDAR hoặc micro thu nhận sóng vật lý và chuyển đổi thành mảng ma trận điểm ảnh số thô $I \in \mathbb{R}^{H \times W \times C}$.
    \item \textbf{Giai đoạn 2: Hiểu (Understanding / Recognition):} Chuyển hóa dữ liệu thô thành ngữ nghĩa, cấu trúc tri thức và mô hình hóa thế giới. Con người hiểu sâu sắc về chức năng, bối cảnh văn hóa và quan hệ nhân quả của thực thể. Trong khi đó, hệ thống AI thực hiện trích xuất đặc trưng phân cấp, ánh xạ xác suất lớp $P(Y \mid X)$, định vị hộp bao và xây dựng đồ thị quan hệ cảnh (Scene Graph).
    \item \textbf{Giai đoạn 3: Thích nghi và Hành động (Action / Adaptation):} Ra quyết định, lập kế hoạch và thực thi hành động để tác động trở lại môi trường hoặc đạt tới mục tiêu tối ưu. Con người hành động dựa trên ý chí, cảm xúc và khả năng ứng biến linh hoạt trước tình huống mới. AI hành động thông qua cơ cấu chấp hành (bộ điều khiển xe tự lái, tín hiệu điều khiển relay) hoặc đề xuất giải pháp, đồng thời thích nghi bằng cách cập nhật trọng số mô hình qua giải thuật học máy.
\end{enumerate}

Từ tiến trình nhận thức trên, bản chất của khái niệm \textbf{Thông minh (Intelligence)} được xác lập dựa trên hai trụ cột cốt lõi:
\begin{itemize}
    \item \textit{Năng lực nhận biết và mô hình hóa vấn đề nhanh chóng:} Khả năng trích xuất thông tin cấu trúc và nắm bắt quy luật tiềm ẩn từ môi trường dữ liệu phức tạp, không đầy đủ hoặc có độ nhiễu cao.
    \item \textit{Năng lực thích ứng linh hoạt (Adaptability):} Khả năng tự điều chỉnh hành vi, cập nhật trọng số tri thức và đưa ra quyết định tối ưu khi môi trường xung quanh hoặc phân phối dữ liệu đầu vào biến đổi.
\end{itemize}

\subsection{Tác động của Trí tuệ Nhân tạo đối với sự phát triển của Trí tuệ Tự nhiên}

Việc ứng dụng Trí tuệ Nhân tạo vào đời sống xã hội và nghiên cứu khoa học mang lại tác động biện chứng hai chiều đối với sự phát triển của trí tuệ con người:

\textbf{1. Khía cạnh tích cực: Khuếch đại năng lực nhận thức (Cognitive Augmentation):}
\begin{itemize}
    \item \textit{Giải phóng con người khỏi các tác vụ tính toán lặp lại:} AI đảm nhận các công việc tốn thời gian nhưng mang tính quy chuẩn cao (như thống kê dữ liệu, rà soát văn bản, phân loại ảnh hàng loạt), tạo điều kiện cho con người tập trung vào tư duy sáng tạo, phản biện và hoạch định chiến lược.
    \item \textit{Mở rộng ranh giới khám phá khoa học:} AI hỗ trợ giải mã các cấu trúc quy luật phức tạp mà trực giác sinh học khó bao quát, điển hình như việc dự đoán cấu trúc protein 3D và các tương tác phân tử sinh học với AlphaFold~\citep{week2_alphafold2021,week2_alphafold2024}, tìm kiếm thiên thể trong thiên văn học hay chẩn đoán sớm tổn thương tế bào y sinh.
\end{itemize}

\textbf{2. Khía cạnh tiêu cực: Nguy cơ phụ thuộc và suy giảm nhận thức (Cognitive Offloading):}
\begin{itemize}
    \item \textit{Suy giảm phản xạ tư duy độc lập:} Sự phụ thuộc quá mức vào các đề xuất tự động của thuật toán (Automation Bias) có thể làm suy yếu khả năng tính toán độc lập, lập luận logic nền tảng và làm giảm sút năng lực ghi nhớ làm việc (Working Memory).
    \item \textit{Hiện tượng định kiến trong bong bóng thông tin (Echo Chambers):} Các thuật toán tối ưu hóa tương tác có thể cô lập người dùng trong các luồng quan điểm một chiều, làm suy giảm tính đa dạng và sự thấu cảm trong nhận thức xã hội.
\end{itemize}

Như vậy, việc ứng dụng AI không tất yếu dẫn đến sự phát triển vượt bậc hay làm thoái hóa trí tuệ tự nhiên. Hệ quả thực tế phụ thuộc vào phương thức con người định vị công nghệ: nếu xem AI là cộng sự khuếch đại nhận thức (Co-pilot) để nâng tầm tư duy thì trí tuệ tự nhiên sẽ tiến bộ; ngược lại nếu sử dụng công nghệ như một giải pháp thay thế thụ động hoàn toàn thì năng lực tư duy tự thân sẽ suy giảm.

\subsection{Phân định bản chất giữa Thông minh (Intelligence) và Thông tuệ (Wisdom)}

Trong khoa học nhận thức và triết học, ``Thông minh'' và ``Thông tuệ'' đại diện cho hai cấp độ nhận thức hoàn toàn khác biệt (Bảng~\ref{tab:comp_intelligence_wisdom}).

\begin{table}[ht]
\centering
\small
\caption{Đối chiếu toàn diện giữa khái niệm Thông minh (Intelligence) và Thông tuệ (Wisdom)}
\label{tab:comp_intelligence_wisdom}
\renewcommand{\arraystretch}{1.3}
\begin{tabularx}{\textwidth}{|p{3.0cm}|X|X|}
\hline
\textbf{Tiêu chí} & \textbf{Thông minh (Intelligence)} & \textbf{Thông tuệ / Trí huệ (Wisdom)} \\
\hline
\textbf{Định nghĩa cốt lõi} & Năng lực giải quyết bài toán và tối ưu hóa mục tiêu: thực thi công việc một cách hiệu quả và chính xác nhất (\textit{Doing things right}). & Năng lực định hướng giá trị và mục đích: lựa chọn đúng mục tiêu có ý nghĩa nhân văn để thực hiện (\textit{Doing the right things}). \\
\hline
\textbf{Đối tượng xử lý} & Dữ liệu số, cấu trúc thông tin, quy tắc logic và các không gian trạng thái toán học. & Ý nghĩa xã hội, chuẩn mực đạo đức, hạnh phúc con người và tính bền vững lâu dài. \\
\hline
\textbf{Động lực vận hành} & Tối ưu hóa hàm mục tiêu (cực tiểu hóa tổn thất hoặc cực đại hóa phần thưởng) dựa trên dữ liệu thống kê quá khứ. & Cân bằng hài hòa giữa lợi ích cá nhân và cộng đồng, gắn liền với lòng trắc ẩn (Compassion) và sự thấu cảm. \\
\hline
\textbf{Khả năng hiện thực} & Máy tính và AI có thể đạt trình độ vượt trội (Super-human) trong các tác vụ tính toán chuyên biệt. & Là đặc tính độc bản của nhận thức sinh học và trải nghiệm sống; AI hiện tại chưa thể có trí tuệ thông tuệ tự thân. \\
\hline
\end{tabularx}
\end{table}

\subsection{Giải thích Luận điểm: ``Thông minh là Khả năng Thích ứng với sự Thay đổi''}

Nhà vật lý lý thuyết Stephen Hawking từng đưa ra định nghĩa nổi tiếng về bản chất của trí tuệ:
\begin{quote}
\textit{``Intelligence is the ability to adapt to change.''} (\textit{Thông minh là khả năng thích ứng với sự thay đổi}.)
\end{quote}

Luận điểm này phản ánh nguyên lý cốt lõi của tính thông minh trong cả tiến hóa sinh học lẫn khoa học máy tính:

\textbf{1. Góc độ Trí tuệ Sinh học và Tâm lý học Nhận thức:}
Môi trường tự nhiên luôn biến động không ngừng. Một sinh vật chỉ duy trì các phản xạ cố định sẽ nhanh chóng bị đào thải khi hoàn cảnh sống thay đổi. Nhà tâm lý học Jean Piaget~\citep{week2_piaget1947} chỉ ra rằng sự phát triển nhận thức thực chất là quá trình thích ứng liên tục thông qua hai cơ chế bổ trợ: \textbf{Đồng hóa (Assimilation)}, tức quá trình tiếp nhận thông tin mới vào cấu trúc tri thức sẵn có; và \textbf{Điều ứng (Accommodation)}, tức việc tái cấu trúc mô hình tư duy để tương thích với thực tế biến đổi. Khả năng điều chỉnh linh hoạt giữa hai cơ chế này giúp sinh vật thích nghi và tồn tại.

\textbf{2. Góc độ Học máy và Thị giác Máy tính:}
Trong kỹ thuật học máy, một mô hình chỉ ghi nhớ máy móc các mẫu trong tập huấn luyện (hiện tượng Overfitting) không thể coi là thông minh. Tính thông minh của hệ thống AI được định lượng bằng năng lực thích ứng với các biến số mới:
\begin{itemize}
    \item \textit{Khả năng khái quát hóa và xử lý dữ liệu lệch phân phối (Out-of-Distribution - OOD):} Mô hình phải nhận diện chính xác đối tượng khi điều kiện chiếu sáng, góc quan sát hoặc môi trường bị biến đổi so với tập dữ liệu huấn luyện ban đầu.
    \item \textit{Thích ứng miền (Domain Adaptation) và Học chuyển giao (Transfer Learning):} Tận dụng tri thức đã học từ miền nguồn (như dữ liệu ảnh mô phỏng đồ họa 3D) để thích ứng nhanh với miền đích thực tế (ảnh camera ngoài đời thực) với lượng dữ liệu gán nhãn tối thiểu.
    \item \textit{Học tăng cường thích ứng (Reinforcement Learning):} Tác tử tự động điều chỉnh chính sách hành động $\pi(a \mid s)$ dựa trên tín hiệu phản hồi phần thưởng từ môi trường động.
    \item \textit{Học trong ngữ cảnh (In-Context Learning):} Khả năng của các Mô hình Nền tảng trong việc thực thi tác vụ mới chỉ qua một vài câu gợi ý hướng dẫn (Prompt) hoặc vài mẫu thị giác mà không cần huấn luyện lại toàn bộ mạng.
\end{itemize}

\subsection{Bản chất của Học máy: Phân tích tư tưởng Arthur Samuel (1959)}

Năm 1959, Arthur Samuel đã đưa ra định nghĩa nền tảng về Học máy:
\begin{quote}
\textit{``Machine Learning: Field of study that gives computers the ability to learn without being explicitly programmed.''}
\end{quote}

Luận điểm trên phân định ranh giới giữa hai hệ hình tính toán:
\begin{itemize}
    \item \textbf{Lập trình truyền thống (Rule-based Systems):} Lập trình viên phải hiểu rõ toàn bộ logic nghiệp vụ, sau đó viết thủ công từng dòng lệnh điều kiện (\texttt{if-else}) và các công thức giải tích. Mô hình vận hành theo sơ đồ:
    \[
    \text{Dữ liệu đầu vào (Data)} + \text{Quy tắc lập trình (Rules)} \longrightarrow \text{Kết quả đầu ra (Output)}
    \]
    Cách tiếp cận này gặp rào cản lớn trước các bài toán thị giác phi cấu trúc, bởi con người không thể lập trình tường minh hàng triệu điều kiện cho mọi biến thiên cường độ pixel của một bức ảnh.
    \item \textbf{Hệ hình Học máy (Learning-based Systems):} Thay vì lập trình quy tắc, kỹ sư thiết kế một kiến trúc mô hình tham số hóa $f_\theta(x)$ và định nghĩa một hàm mục tiêu (Loss Function) $\mathcal{L}$. Thuật toán sẽ tự động tối ưu hóa tập tham số $\theta$ dựa trên dữ liệu mẫu:
    \[
    \text{Dữ liệu đầu vào (Data)} + \text{Dữ liệu nhãn/Mục tiêu (Output)} \longrightarrow \text{Mô hình đã học ($\theta^*$)}
    \]
\end{itemize}

\subsection{Định nghĩa Hình thức về Học máy theo Tom Mitchell (1997)}

Năm 1997, Tom Mitchell đã chuẩn hóa lý thuyết học máy bằng định nghĩa toán học chặt chẽ:
\begin{quote}
\textit{``A computer program is said to learn from experience $\mathcal{E}$ with respect to some class of tasks $\mathcal{T}$ and performance measure $\mathcal{P}$, if its performance at tasks in $\mathcal{T}$, as measured by $\mathcal{P}$, improves with experience $\mathcal{E}$.''}
\end{quote}

Định nghĩa này cấu thành từ bộ ba $(\mathcal{T}, \mathcal{E}, \mathcal{P})$:
\begin{itemize}
    \item \textbf{Nhiệm vụ $\mathcal{T}$ (Task):} Bài toán mục tiêu mà hệ thống cần thực hiện. Ví dụ trong thị giác máy tính: Phát hiện vị trí các khuyết tật trên bề mặt vi mạch $I \in \mathbb{R}^{H \times W \times C}$.
    \item \textbf{Trải nghiệm $\mathcal{E}$ (Experience):} Tập hợp dữ liệu lịch sử hoặc các tương tác mà mô hình được tiếp cận. Ví dụ: Tập dữ liệu $\mathcal{D} = \{(x_i, y_i)\}_{i=1}^N$ gồm $100,000$ hình ảnh vi mạch được gán nhãn vị trí lỗi chuẩn bởi chuyên gia kiểm thử.
    \item \textbf{Độ đo hiệu năng $\mathcal{P}$ (Performance Measure):} Thước đo định lượng toán học dùng để đánh giá chất lượng thực thi tác vụ. Ví dụ: Chỉ số độ chính xác trung bình $\text{mAP}@0.5$ trên tập kiểm thử độc lập $\mathcal{D}_{\text{test}}$.
\end{itemize}

Một thuật toán được xác nhận là đang học khi chất lượng thực thi tăng đơn điệu theo lượng dữ liệu trải nghiệm:
\begin{equation}
\frac{\partial \mathcal{P}(\mathcal{T})}{\partial \mathcal{E}} > 0 \quad \text{hoặc} \quad \lim_{|\mathcal{E}| \to \infty} \mathcal{L}(\theta; \mathcal{E}) \to \min
\end{equation}

\subsection{Tiến trình Chuyển tiếp từ Học máy sang Học sâu: Giới hạn và Động lực Đột phá}

Sự phát triển của thị giác máy tính gắn liền với bước chuyển dịch từ các thuật toán học máy cổ điển sang các kiến trúc mạng nơ-ron sâu:

\textbf{1. Quá trình chuyển dịch công nghệ:}
Trong thời kỳ học máy truyền thống (trước năm 2012), quy trình xử lý tách rời làm hai giai đoạn độc lập: (1) Trích xuất đặc trưng thủ công (Hand-crafted Feature Extraction) do con người thiết kế như SIFT~\citep{lowe2004distinctive}, HOG~\citep{dalal2005histograms} hay LBP; (2) Đưa các vector đặc trưng này vào bộ phân loại nông như SVM~\citep{cortes1995support} hoặc Random Forest~\citep{breiman2001random}. Ngược lại, kỷ nguyên học sâu tích hợp toàn bộ quy trình thành một mạng duy nhất theo cơ chế học biểu diễn đầu-cuối (End-to-End Representation Learning), ánh xạ trực tiếp từ pixel thô sang dự đoán ngữ nghĩa.

\textbf{2. Những hạn chế căn bản của Học máy truyền thống:}
\begin{itemize}
    \item \textit{Nút thắt cổ chai từ đặc trưng thủ công:} Các bộ trích xuất thủ công chỉ nắm bắt được các mẫu hình học cục bộ đơn giản (cạnh, góc, kết cấu xám), dễ mất tính bất biến trước các biến dạng phối cảnh lớn, che khuất hoặc thay đổi tư thế phức tạp.
    \item \textit{Lời nguyền số chiều (Curse of Dimensionality):} Khi xử lý ảnh có độ phân giải cao, không gian đặc trưng bùng nổ khiến các thuật toán dựa trên ranh giới phẳng hoặc khoảng cách hình học bị suy giảm hiệu năng nghiêm trọng.
    \item \textit{Sự bão hòa dung lượng mô hình:} Khi tăng quy mô dữ liệu từ hàng vạn lên hàng triệu mẫu, hiệu năng của các mô hình nông nhanh chóng đạt trần do dung lượng biểu diễn (Model Capacity) bị giới hạn.
\end{itemize}

\textbf{3. Động lực thúc đẩy sự ra đời của Học sâu:}
\begin{itemize}
    \item \textit{Tự động học biểu diễn phân cấp (Hierarchical Feature Abstraction):} Mạng nơ-ron sâu tự động trừu tượng hóa thông tin theo từng tầng: từ đường biên cục bộ ở tầng đầu, đến các bộ phận hình thể ở tầng giữa, và ngữ nghĩa đối tượng hoàn chỉnh ở tầng sâu nhất.
    \item \textit{Định luật mở rộng quy mô (Scaling Laws):} Hiệu năng của học sâu liên tục tăng trưởng tương ứng với kích thước dữ liệu và số lượng tham số khi được hỗ trợ bởi năng lực tính toán GPU/TPU song song quy mô lớn.
    \item \textit{Năng lực biểu diễn đa phương thức:} Khả năng liên kết biểu diễn đồng thời ảnh, video, văn bản và âm thanh vào cùng một không gian nhúng chung.
\end{itemize}

\subsection{Nền tảng của Học sâu: Vì sao Nguyên lý Đơn giản tạo nên Khả năng Biểu diễn Vượt trội?}

Về mặt toán học, mạng nơ-ron sâu được xây dựng từ các phép toán nền tảng:
\begin{enumerate}
    \item Phép biến đổi affine tuyến tính: $\mathbf{z} = \mathbf{W}\mathbf{x} + \mathbf{b}$.
    \item Phép kích hoạt phi tuyến cục bộ từng phần tử: $\mathbf{a} = \sigma(\mathbf{z})$ (với $\sigma$ là hàm ReLU, GELU, Sigmoid).
    \item Quá trình tối ưu hóa bằng giải thuật lan truyền ngược (Backpropagation) dựa trên quy tắc đạo hàm chuỗi (Chain Rule) và phương pháp hạ dốc độ dốc ngẫu nhiên (SGD~\citep{robbins1951stochastic}).
\end{enumerate}

Nguyên lý toán học này mang lại khả năng biểu diễn vượt trội nhờ ba cơ chế nền tảng:
\begin{itemize}
    \item \textbf{Định lý Xấp xỉ Phổ quát (Universal Approximation Theorem):} Một mạng nơ-ron chứa các tầng ẩn phi tuyến có khả năng xấp xỉ bất kỳ hàm số liên tục nào với độ chính xác tùy ý. Khi xếp chồng nhiều tầng sâu, cấu trúc phân cấp cho phép mô hình biểu diễn các hàm logic phức tạp với số lượng tham số tiết kiệm hơn theo cấp số nhân so với mạng nông.
    \item \textbf{Cấu trúc thiên lệch cảm ứng (Inductive Bias):} Kiến trúc tích chập (Convolution) tận dụng tính chất bất biến tịnh tiến và tính tương quan cục bộ của hình ảnh tự nhiên, giúp thu hẹp không gian tìm kiếm tham số một cách tối ưu. Tương tự, cơ chế tự chú ý (Self-Attention) trong Vision Transformer mô hình hóa quan hệ phụ thuộc toàn cục linh hoạt.
    \item \textbf{Hình học tối ưu trong không gian tham số cao chiều:} Trong không gian hàng triệu đến hàng tỷ chiều, các điểm dừng tối ưu chủ yếu là các điểm yên ngựa (Saddle Points) mà giải thuật tối ưu hóa có thể vượt qua, giúp mô hình hội tụ về các vùng nghiệm có khả năng khái quát hóa cao.
\end{itemize}

\subsection{Mô hình 5 Cấp độ Nhận thức và Dữ liệu: Ứng dụng Giám sát An ninh Trường học}

Quá trình chuyển hóa dữ liệu từ dạng tín hiệu vật lý thô sang hành động có giá trị nhân văn được mô hình hóa qua \textbf{5 Cấp độ Nhận thức}:
\begin{enumerate}
    \item \textbf{Level 1: Unfiltered Data (Dữ liệu thô chưa lọc):} Tín hiệu điểm ảnh, sóng âm, dữ liệu cảm biến chưa qua xử lý.
    \item \textbf{Level 2: Pattern Information (Thông tin mẫu):} Các mẫu hình học, đặc trưng thị giác, vector nhúng được trích xuất có cấu trúc.
    \item \textbf{Level 3: Predictability Knowledge (Tri thức có tính dự đoán):} Nhận định danh tính, xu hướng, mối liên hệ nhân quả và quy luật lịch sử.
    \item \textbf{Level 4: Choice Intelligence (Trí thông minh ra quyết định):} Năng lực lựa chọn phương án hành động tối ưu trong các kịch bản thực thi.
    \item \textbf{Level 5: Compassion Wisdom (Thông tuệ thấu cảm):} Sự đánh giá giá trị nhân văn, thấu hiểu hoàn cảnh và đưa ra giải pháp nhân bản.
\end{enumerate}

Bảng~\ref{tab:matrix_5_levels_school} phân tích ma trận ứng dụng 5 cấp độ này vào bài toán \textbf{Hệ thống Giám sát An ninh \& Điểm danh Khuôn mặt tại Trường học}, chỉ rõ từng Task, Cấp độ nhận thức tương ứng, Dữ liệu đầu vào, Dữ liệu đầu ra và Nội dung chi tiết trong từng ô.

\begin{table}[ht]
\centering
\scriptsize
\setlength{\tabcolsep}{3.5pt}
\caption{Ma trận phân tích 5 cấp độ nhận thức ứng dụng trong Hệ thống Giám sát An ninh Trường học}
\label{tab:matrix_5_levels_school}
\renewcommand{\arraystretch}{1.3}
\begin{tabularx}{\textwidth}{|p{2.1cm}|p{2.0cm}|p{2.5cm}|p{2.4cm}|X|}
\hline
\textbf{Task hệ thống} & \textbf{Cấp độ nhận thức} & \textbf{Dữ liệu đầu vào (Input)} & \textbf{Dữ liệu đầu ra (Output)} & \textbf{Nội dung chi tiết \& Ví dụ minh họa trong ô} \\
\hline
\textbf{Task 1: Thu nhận luồng hình ảnh} & \textbf{Level 1: Unfiltered Data} & Luồng tín hiệu video RTSP từ camera giám sát tại cổng trường & Mảng ma trận byte điểm ảnh thô $I_t \in \mathbb{R}^{H \times W \times 3}$ & \textbf{[Mức 1 - Dữ liệu thô]:} Thu nhận 30 khung hình/giây chứa mọi thông tin nền (bóng râm, cây cối, xe cộ, ánh sáng đổi) mà máy tính chưa phân biệt được có học sinh hay không. \\
\hline
\textbf{Task 2: Phát hiện \& Trích xuất khuôn mặt} & \textbf{Level 2: Pattern Information} & Khung ảnh thô $I_t$ & Hộp bao khuôn mặt $\mathcal{B}$ và Vector nhúng đặc trưng $\mathbf{e} \in \mathbb{R}^{512}$ & \textbf{[Mức 2 - Thông tin mẫu]:} Thuật toán RetinaFace định vị khuôn mặt và trích xuất vector 512 chiều biểu diễn cấu trúc hình học khoảng cách mắt, mũi, miệng. \\
\hline
\textbf{Task 3: Định danh \& Đối chiếu lịch sử} & \textbf{Level 3: Predictability Knowledge} & Vector nhúng $\mathbf{e}$ kèm cơ sở dữ liệu học sinh $\mathcal{D}_{\text{student}}$ & Danh tính học sinh kèm độ tin cậy, lịch sử chuyên cần & \textbf{[Mức 3 - Tri thức dự đoán]:} Xác định khuôn mặt là ``Em An - Lớp 10A1'' với độ khớp 99.4\%; nhận diện học sinh đi trễ 3 ngày liên tiếp, dự báo nguy cơ vi phạm nền nếp. \\
\hline
\textbf{Task 4: Điều khiển cổng \& Cảnh báo tự động} & \textbf{Level 4: Choice Intelligence} & Kết quả định danh và quy tắc an ninh trường học & Tín hiệu relay mở cổng / Gửi thông báo phụ huynh / Chuông báo & \textbf{[Mức 4 - Trí thông minh lựa chọn]:} Nếu là học sinh: tự động mở chốt cửa và gửi thông báo điểm danh cho phụ huynh; nếu phát hiện người lạ chưa đăng ký: tự động khóa cổng và báo bảo vệ. \\
\hline
\textbf{Task 5: Đánh giá hoàn cảnh \& Hỗ trợ nhân văn} & \textbf{Level 5: Compassion Wisdom} & Nhật ký đi trễ, biểu cảm khuôn mặt mệt mỏi và dữ liệu gia cảnh & Đề xuất tư vấn tâm lý, hỗ trợ hoàn cảnh khó khăn & \textbf{[Mức 5 - Thông tuệ thấu cảm]:} Không trừ điểm thi đua máy móc; hệ thống ghi nhận học sinh đi trễ kèm biểu hiện kiệt sức, đề xuất giáo viên chủ nhiệm thăm hỏi gia cảnh và hỗ trợ tâm lý học đường. \\
\hline
\end{tabularx}
\end{table}

\subsection{Mô hình Giải pháp 3 Tầng trong Khoa học Dữ liệu Thị giác}

Để xây dựng một hệ thống thị giác máy tính hoàn chỉnh phục vụ thực tiễn, kiến trúc kỹ thuật được chuẩn hóa thành \textbf{Mô hình Giải pháp 3 Tầng} (Bảng~\ref{tab:three_tier_cv_architecture}):
\begin{enumerate}
    \item \textbf{Tầng 1: Kỹ thuật Cốt lõi / Nền tảng Thuật toán (Foundational Technologies):} Bao gồm các mô hình toán học, kiến trúc mạng nơ-ron nền tảng (CNN, Vision Transformer), hàm mất mát, thuật toán tối ưu (AdamW~\citep{loshchilov2019decoupled}, SGD~\citep{robbins1951stochastic}) và giải thuật hình học thị giác (Epipolar, Homography).
    \item \textbf{Tầng 2: Tác vụ Thị giác Cơ sở (Core Vision Tasks):} Các khối chức năng thị giác cơ bản giải quyết các bài toán độc lập như Phân loại ảnh, Phát hiện hộp bao, Nhận dạng khuôn mặt, Ước lượng chiều sâu, Tách mặt nạ phân đoạn hay Theo dõi quỹ đạo.
    \item \textbf{Tầng 3: Hệ thống Ứng dụng Nghiệp vụ Hoàn chỉnh (End-to-End Business Applications):} Tích hợp các khối chức năng ở Tầng 2 cùng các ràng buộc nghiệp vụ, giao diện người dùng, cơ sở dữ liệu và phần cứng điều khiển để tạo thành sản phẩm hoàn chỉnh giải quyết bài toán thực tế.
\end{enumerate}

\begin{table}[ht]
\centering
\small
\caption{Cấu trúc Mô hình Giải pháp 3 Tầng trong Thị giác máy tính}
\label{tab:three_tier_cv_architecture}
\renewcommand{\arraystretch}{1.3}
\begin{tabularx}{\textwidth}{|p{2.8cm}|X|X|}
\hline
\textbf{Phân tầng} & \textbf{Thành phần cốt lõi} & \textbf{Ví dụ minh họa cụ thể} \\
\hline
\textbf{Tầng 1: Kỹ thuật cốt lõi} & Kiến trúc mạng (ResNet~\citep{he2016deep}, Swin-ViT~\citep{liu2021swin}, YOLO-backbone~\citep{redmon2016you}), Hàm mất mát (ArcFace, InfoNCE, GIoU~\citep{rezatofighi2019generalized}), Thư viện tối ưu (PyTorch, TensorRT). & Mạng ResNet-50~\citep{he2016deep} kết hợp hàm mất mát ArcFace tối ưu hóa bằng AdamW~\citep{loshchilov2019decoupled} trên tập dữ liệu chuẩn. \\
\hline
\textbf{Tầng 2: Tác vụ thị giác} & Các bài toán cơ sở: Face Detection, Facial Landmark Alignment, Re-Identification, Crowd Density Estimation. & Phát hiện 68 điểm mốc khuôn mặt và trích xuất vector đặc trưng khuôn mặt 512 chiều. \\
\hline
\textbf{Tầng 3: Ứng dụng nghiệp vụ} & Hệ thống hoàn chỉnh tích hợp cơ sở dữ liệu, phân quyền, cảnh báo thời gian thực và tự động hóa quy trình nghiệp vụ. & Hệ thống Kiosk tự động Check-in hành khách tại sân bay tích hợp cổng kiểm soát an ninh xuất nhập cảnh. \\
\hline
\end{tabularx}
\end{table}

\subsection{Khái quát các Nhóm Ứng dụng Cơ bản của Thị giác Máy tính}

Thị giác máy tính hiện đại được phân bổ thành 5 nhóm ứng dụng công nghiệp trọng yếu (Bảng~\ref{tab:cv_application_domains}).

\begin{table}[ht]
\centering
\small
\caption{Tổng hợp 5 nhóm lĩnh vực ứng dụng cơ bản của Thị giác máy tính}
\label{tab:cv_application_domains}
\renewcommand{\arraystretch}{1.3}
\begin{tabularx}{\textwidth}{|p{3.2cm}|X|X|}
\hline
\textbf{Lĩnh vực ứng dụng} & \textbf{Các tác vụ Tầng 2 tích hợp} & \textbf{Hệ thống sản phẩm Tầng 3 tiêu biểu} \\
\hline
\textbf{1. Giám sát an ninh \& Đô thị thông minh} & Detection, Tracking, License Plate Recognition, Crowd Counting. & Hệ thống camera phạt nguội giao thông tự động; Giám sát trật tự đô thị và an ninh sân bay. \\
\hline
\textbf{2. Y tế số \& Chẩn đoán hình ảnh} & Semantic Segmentation, 3D Voxel Reconstruction, Anomaly Detection. & Phần mềm tự động khoanh vùng khối u trên ảnh chụp CT/MRI; Phát hiện sớm polyp nội soi tiêu hóa. \\
\hline
\textbf{3. Tự hành \& Robot công nghiệp} & 3D Object Detection, Depth Estimation, Visual SLAM, Optical Flow. & Hệ thống xe tự lái; Cánh tay robot tự động gắp thả linh kiện và kiểm tra lỗi hàn vi mạch. \\
\hline
\textbf{4. Bán lẻ \& Thương mại điện tử} & Image Retrieval, Fine-grained Classification, Pose Estimation. & Tìm kiếm sản phẩm bằng hình ảnh trên sàn TMĐT; Cửa hàng bán lẻ tự động không người thu ngân. \\
\hline
\textbf{5. Điện ảnh, CGI \& Thực tế ảo (AR/VR)} & 3D Face Reconstruction, Motion Capture, Generative Video Inpainting. & Kỹ xảo điện ảnh kỹ thuật số (VFX); Tạo Avatar Metahuman; Bộ lọc hiệu ứng khuôn mặt AR trên smartphone. \\
\hline
\end{tabularx}
\end{table}

\subsection{Phân tích Chuyên sâu: Ứng dụng Trí tuệ Nhân tạo trong Ngành Điện ảnh}

Ứng dụng AI trong phân tích trải nghiệm khán giả điện ảnh (AI for Cinema Audience Analysis) là ví dụ điển hình về sự phối hợp giữa Tầng 1, Tầng 2 và Tầng 3 nhằm tối ưu hóa giá trị nghệ thuật và thương mại (Bảng~\ref{tab:ai_cinema_pipeline}):

\begin{table}[ht]
\centering
\small
\caption{Quy trình phân tích trải nghiệm khán giả điện ảnh ứng dụng AI}
\label{tab:ai_cinema_pipeline}
\renewcommand{\arraystretch}{1.3}
\begin{tabularx}{\textwidth}{|p{3.0cm}|X|X|}
\hline
\textbf{Giai đoạn phân tích} & \textbf{Công nghệ \& Tác vụ thị giác (Tầng 1 \& 2)} & \textbf{Giá trị ứng dụng nghiệp vụ điện ảnh (Tầng 3)} \\
\hline
\textbf{1. Nhận diện khuôn mặt \& Cảm xúc} & Camera hồng ngoại trong rạp tối $\rightarrow$ Facial Landmark Detection $\rightarrow$ Phân loại biểu cảm (Vui mừng, Sợ hãi, Bất ngờ, Thất vọng) theo hệ thống mã hóa hành vi nét mặt FACS. & Đo lường chính xác mức độ kích thích cảm xúc của người xem tại từng cảnh phim; kiểm chứng xem phân đoạn phim có đạt hiệu ứng mong muốn (ví dụ cảnh hài có làm khán giả cười hay không). \\
\hline
\textbf{2. Phân tích hành vi \& Hướng nhìn} & Head Pose Estimation, Eye Gaze Tracking $\rightarrow$ Action Recognition (che mắt, ngả người ra sau, quay sang người bên cạnh, xem điện thoại). & Đánh giá mức độ tập trung (Engagement Level) của khán giả; phát hiện các phân đoạn gây buồn ngủ hoặc nhàm chán khiến khán giả phân tâm. \\
\hline
\textbf{3. Đồng bộ Timeline \& Tối ưu kịch bản} & Tích hợp chuỗi thời gian cảm xúc khán giả đồng bộ với biểu đồ Timeline dựng phim (Cut-by-cut synchronization). & Hỗ trợ đạo diễn và nhà sản xuất trong khâu chiếu thử nghiệm (Test Screening): cắt gọt các phân đoạn lê thê, tinh chỉnh nhịp độ (Pacing), dự báo doanh thu phòng vé và định hướng chiến dịch Marketing. \\
\hline
\end{tabularx}
\end{table}

\subsection{Ý nghĩa Thực tiễn của Bài toán Đếm Đám đông ở Tầng Ứng dụng Nghiệp vụ}

Ở Tầng 2 (Tác vụ kỹ thuật), bài toán Đếm đám đông (Crowd Counting) giải quyết việc ánh xạ hình ảnh $I \in \mathbb{R}^{H \times W \times C}$ sang bản đồ mật độ liên tục $D(x, y)$ để ước lượng số thực thể $\hat{N} = \sum_{x,y} D(x, y)$.

Khi được triển khai lên \textbf{Tầng 3 (Hệ thống Ứng dụng Nghiệp vụ Thực tiễn)}, bài toán này mang lại giá trị quản trị và an toàn xã hội to lớn:
\begin{itemize}
    \item \textbf{Phòng ngừa thảm họa giẫm đạp (Stampede Prevention):} Tại các sự kiện lễ hội, nhà ga tàu điện ngầm hoặc sân vận động, hệ thống liên tục tính toán mật độ cục bộ ($\text{người}/\text{m}^2$). Khi mật độ vượt ngưỡng an toàn ($>4\text{ người}/\text{m}^2$), hệ thống tự động kích hoạt cảnh báo nguy cơ tắc nghẽn cục bộ để lực lượng an ninh kịp thời phân luồng trước khi xảy ra tai nạn.
    \item \textbf{Tối ưu hóa phân bổ nguồn lực công cộng:} Dữ liệu mật độ thời gian thực giúp cơ quan quản lý điều phối tần suất phương tiện giao thông công cộng, bố trí nhân viên cứu hộ và mở thêm cửa thoát hiểm linh hoạt.
    \item \textbf{Phân tích hành vi tiêu dùng và quy hoạch thương mại:} Đếm lưu lượng người đi bộ qua các gian hàng trong trung tâm thương mại giúp xác định các vị trí tập trung đông khách để định giá mặt bằng cho thuê và tối ưu hóa vị trí đặt biển quảng cáo.
\end{itemize}

\subsection{Phân tích Hệ thống Tái tạo Khuôn mặt 3D theo Mô hình Giải pháp 3 Tầng}

Hệ thống tái tạo khuôn mặt 3D (3D Face Reconstruction) từ ảnh 2D đơn góc nhìn là minh chứng tiêu biểu cho kiến trúc phân tầng hoàn chỉnh từ lý thuyết toán học đến ứng dụng thực tế (Bảng~\ref{tab:three_tier_3d_face}).

\begin{table}[ht]
\centering
\small
\caption{Phân tích hệ thống Tái tạo Khuôn mặt 3D theo Mô hình Giải pháp 3 Tầng}
\label{tab:three_tier_3d_face}
\renewcommand{\arraystretch}{1.3}
\begin{tabularx}{\textwidth}{|p{2.8cm}|X|X|}
\hline
\textbf{Phân tầng} & \textbf{Nhiệm vụ \& Thành phần kỹ thuật} & \textbf{Chi tiết giải pháp công nghệ} \\
\hline
\textbf{Tầng 1: Kỹ thuật cốt lõi} & Cơ sở mô hình toán hình học 3D, kiến trúc mạng nơ-ron và hàm tối ưu khả vi (Differentiable Rendering). & Mô hình tham số hình học 3D Morphable Model (3DMM: BFM, FLAME); Mạng Transformer trích xuất hệ số hình dạng $\boldsymbol{\alpha}$, biểu cảm $\boldsymbol{\beta}$ và tư thế $\mathbf{p}$; Hàm mất mát Photometric Consistency Loss, Landmark Reprojection Loss, và Perceptual Identity Loss. \\
\hline
\textbf{Tầng 2: Tác vụ thị giác cơ sở} & Các khối trích xuất thông tin trung gian từ ảnh 2D sang không gian 3D. & Phát hiện 68/468 điểm mốc 3D (Dense Facial Landmarks); Ước lượng trường độ sâu đơn sắc (Monocular Depth Estimation); Tách bản đồ phản xạ ánh sáng (Albedo/Texture UV Map) và góc chiếu sáng (Spherical Harmonics). \\
\hline
\textbf{Tầng 3: Ứng dụng nghiệp vụ hoàn chỉnh} & Tích hợp vào các phần mềm ứng dụng phục vụ người dùng cuối. & \begin{itemize}[leftmargin=*,noitemsep,topsep=0pt]
    \item \textbf{Y tế \& Thẩm mỹ:} Mô phỏng kết quả phẫu thuật chỉnh hình hàm mặt trước khi phẫu thuật thực tế.
    \item \textbf{Giải trí \& Game:} Tự động tạo Digital Avatar 3D cá nhân hóa có khả năng cử động khớp thời gian thực trong game và không gian thực tế ảo.
    \item \textbf{Bảo mật sinh trắc học:} Xác thực khuôn mặt 3D chống giả mạo (Face Anti-spoofing) ngăn chặn hành vi vượt rào bằng ảnh in 2D hoặc màn hình kỹ thuật số.
\end{itemize} \\
\hline
\end{tabularx}
\end{table}
